\chapter{Συμπεράσματα και μελλοντική εργασία}
\label{ch:conclusions}

\guide{Έκταση 3--4 σελίδες. Καμία νέα πληροφορία· σύνοψη των συνεισφορών με τα κύρια αποτελέσματα.}

\section{Συμπεράσματα}
\label{sec:conclusions}
\todo{Σύνοψη.}

\section{Μελλοντική εργασία}
\label{sec:future}
\begin{itemize}
  \item \textbf{Αυτόματη επαναευθυγράμμιση} με ένα anchor ανά σημείο ελέγχου· η \cref{sec:exp-anchors} είναι η «Φάση~0» αυτής της πρότασης (\code{docs/proposals/auto-realignment.md}).
  \item \textbf{Διανύσματα μετακίνησης 3Δ και σύγκριση δύο ημερομηνιών}, με ένα endpoint «στιγμιότυπο έργου» στην πλατφόρμα (\code{docs/proposals/terraweb-data-for-new-features.md}, \code{displacement-visualization.md}).
  \item \textbf{Ευθυγράμμιση χωρίς σημεία ελέγχου}, με ταίριασμα του πλέγματος του χώρου που καταγράφει η συσκευή με το μοντέλο φωτογραμμετρίας (ICP \cite{besl1992icp}).
  \item \textbf{Συνεργατική χρήση}: κοινά anchors για πολλούς χρήστες στο ίδιο εργοτάξιο.
  \item \textbf{Φορητότητα}: OpenXR και Universal Render Pipeline για άλλα headsets.
\end{itemize}
